\chapter{Mpox}
\index{Mpox}
\label{sec:Mpox}

\section{Overview}
Mpox is a viral zoonosis caused by mpox virus, a double-stranded DNA orthopoxvirus closely related to variola virus, formerly called monkeypox, endemic in Central and West Africa, with the 2022 global outbreak involving over 90,000 cases across 110 countries.
\section{Causes}
This condition happens because of the virus entering through broken skin, respiratory tract, or mucous membranes, replicating at the inoculation site, spreading to regional lymph nodes, and causing viremia.     
\section{Risk Factors}

\begin{itemize}[noitemsep]
  \item Genetic predisposition and family history
  \item Advancing age
  \item Lifestyle factors (diet, exercise, smoking, alcohol)
  \item Environmental exposures
  \item Underlying medical conditions
  \item Medication use
\end{itemize}
\section{Signs and Symptoms}

\subsection{Common symptoms}
\begin{itemize}[noitemsep]
  \item \item You may experience an incubation of 5--21 days, with prodrome (fever, headache, muscle pain, lymphadenopathy---hallmark distinguishing mpox from smallpox) followed by rash appearing 1--4 days after fever, progressing from macules to papules, vesicles, pustules, and crusts over 2--4 weeks, with a centrifugal distribution, and anogenital lesions common in the 2022 outbreak. Pain or discomfort in the affected area    \item Pain or discomfort in the affected area
  \end{itemize}

\subsection{Emergency warning signs}
\begin{itemize}[noitemsep]
  \item Severe or worsening symptoms of mpox require immediate medical evaluation
\end{itemize}
\section{Progression and Stages}
The condition may progress through different stages of severity over time. Early stages often involve mild or subtle symptoms that may be overlooked. Moderate stages involve more noticeable symptoms affecting daily function. Advanced stages may involve significant impairment and complications.
\section{Complications}
\begin{itemize}[noitemsep]
  \item Disease progression leading to worsening symptoms
  \item Secondary infections or injuries
  \item Side effects of treatment
  \item Reduced quality of life
  \item Emotional and psychological impact
\end{itemize}
\section{Diagnosis}
Doctors diagnose this using PCR of lesion swabs. Comprehensive medical history and physical examination Laboratory tests to assess organ function and rule out other conditions Imaging studies to visualize affected structures Specialized testing depending on the affected system Consultation with relevant specialists Monitoring of disease progression over time

\begin{itemize}[noitemsep]
  \item Comprehensive medical history and physical examination
  \item Laboratory tests to assess organ function
  \item Imaging studies to visualize affected structures
  \item Specialized testing depending on the affected system
  \item Consultation with relevant specialists
\end{itemize}
\section{Prevention}
\begin{itemize}[noitemsep]
  \item Maintain a healthy lifestyle with balanced diet and exercise
  \item Avoid known risk factors
  \item Attend regular medical check-ups for early detection
  \item Follow recommended screening guidelines
\end{itemize}
\section{Treatment Options}
Medications to manage symptoms and address underlying mechanisms Lifestyle modifications and self-care measures Physical therapy and rehabilitation Surgical intervention when indicated Regular monitoring and follow-up care Patient education and support services

\begin{itemize}[noitemsep]
  \item Medications to manage symptoms and address underlying mechanisms
  \item Lifestyle modifications and self-care measures
  \item Physical therapy and rehabilitation
  \item Surgical intervention when indicated
  \item Regular monitoring and follow-up care
\end{itemize}
\section{Living with It}
Treatment options include tecovirimat (TPOXX) and supportive care.

\begin{itemize}[noitemsep]
  \item Follow your treatment plan as prescribed
  \item Maintain regular follow-up with your healthcare provider
  \item Make necessary lifestyle adjustments
  \item Seek support from family, friends, or support groups
  \item Stay informed about your condition
\end{itemize}
\section{Outlook}
\begin{itemize}[noitemsep]
  \item In most cases, the outlook is good
  \item Clade I (Central Africa) is more virulent than clade II (West Africa).
\end{itemize}
